% Ad Familiares 13.18 (ad-familiares-13-18) --- English translation
% Translated by Claude (Anthropic) from the Latin (Perseus).
% Style guide: STYLE.md.

\ciceroLetterOpener{CICERO SERVIO S.}

\ciceroSection{1} I will not grant that the letter you wrote to our Atticus, so charming in itself and so very civilly composed, was more agreeable to him --- whom I saw beside himself with delight --- than it was to me. For although it was almost equally welcome to both of us, still I admired you the more on this account, that, where a generous reply from you would have come naturally if you had been asked or at any rate prompted (and even so, there was no doubt in our minds that such would have been your response), you went further: you wrote to him unsolicited, and conveyed to him by letter, when he was not looking for it, so signal a mark of your good will. On that score, I am not bound to ask you to act with the more enthusiasm on my behalf as well --- for nothing fuller than what you are promising could be done --- nor even to thank you for what you have done both for his sake and of your own accord;

\ciceroSection{2} but I will say this much: what you have done is most welcome to me. For such a judgment of yours about a man whom I, above all others, hold dear cannot but be in the highest degree agreeable to me; and being so, it cannot fail to be grounds for gratitude. And yet, since the closeness of our friendship gives me leave to transgress with you even in the matter of writing, I shall do both the things I just said I should not. For of what you have shown yourself ready to do on Atticus's account, I would have you add as much more as can be added on the strength of our affection; and where just now I was hesitating to thank you, I now thank you outright, and I would have you so understand it: by whatever good offices, in the affairs of Epirus and the rest, you bind Atticus to yourself, by those same offices you will lay me under obligation to you.
